\documentclass[10pt,a5paper]{article}
\usepackage[utf8]{inputenc}
\usepackage[T1]{fontenc}
\usepackage[ngerman]{babel}
\usepackage{geometry}
\geometry{a5paper, margin=1.2cm}
\usepackage{amsmath, amssymb}
\usepackage{microtype}
\usepackage{hyperref}
\pagestyle{empty}

\begin{document}

\begin{center}
    {\Huge\bfseries WDBT+ rettet die Welt!}

    \vspace{0.2cm}
    {\large\itshape Die frohe Botschaft der Weber-de Broglie-Bohm-Theorie}
    
    \vspace{0.1cm}
    {\small (plus fraktale Dimension – weil's sonst zu einfach wäre)}
\end{center}

\vspace{0.3cm}
\hrule
\vspace{0.3cm}

\subsection*{Was war nochmal das Problem?}
Die alte Physik: 
\begin{itemize}
    \item ART: Raumzeit krümmt sich (warum? \textit{»Weil's so ist.«})
    \item Quantenmechanik: Zufall! (\textit{»Wir würfeln, weil's keiner checkt.«})
    \item Dunkle Materie: \textit{»Da muss was sein – finden tun wir's nicht.«}
    \item Dunkle Energie: \textit{»Äh, da ist noch mehr – keine Ahnung was.«}
\end{itemize}

\subsection*{Die WDBT+ Lösung:}
\textbf{Ein Framework, sie alle zu knechten, ins Dunkle zu treiben und ewig zu binden.}

\subsubsection*{1. Weber-Elektrodynamik (WED)}
\[
F = \frac{q_1q_2}{4\pi\epsilon_0 r^2}\left[1 - \frac{\dot{r}^2}{c^2} + 2\frac{r\ddot{r}}{c^2}\right]
\]
\textit{»Maxwell ist nur der langweilige Mittelwert.«}

\subsubsection*{2. Weber-Gravitation (WG)}
\[
F = -\frac{GMm}{r^2}\left[1 - \frac{\dot{r}^2}{c^2} + 0,5\frac{r\ddot{r}}{c^2}\right]
\]
\textit{»Periheldrehung? Krieg' ich hin. Ohne Raumzeit-Krümmung. Bitte sehr.«}

\subsubsection*{3. Quantenpotential (DBT)}
\[
Q = -\frac{\hbar^2}{2m}\frac{\nabla^2\sqrt{\rho}}{\sqrt{\rho}}
\]
\textit{»Schrödinger ist nur die Hülle. Die Führung kommt von Q. Deterministisch, Baby!«}

\subsubsection*{4. Fraktale Dimension}
\[
D = \frac{\ln 20}{\ln(2+\phi)} \approx 2,71
\]
\textit{»3 ist so 20. Jahrhundert. Wir sind bei 2,71 – der Goldene Schnitt lässt grüßen.«}

\subsection*{Was folgt daraus? (Achtung, Mind-blowing!)}

\begin{itemize}
    \item \textbf{Rotation ist fundamental:} Newton setzt sie voraus, wir \textit{erzeugen} sie. Neidisch?
    \item \textbf{Hot Jupiters:} Einfach die Einwärtsphase. Keine komplizierte Migration.
    \item \textbf{Monddrift:} 3,8 cm/Jahr – Gezeiten? Nein, DSTT! (ART ist nur die kleine Schwester)
    \item \textbf{Galaxienrotation:} Ohne Dunkle Materie. Einfach fraktale Struktur. Zack, flach.
    \item \textbf{Lichtablenkung:} ART sagt: alle Frequenzen gleich. WDBT+ sagt: \textit{»Nö, kommt auf die Farbe an.«} Testen!
\end{itemize}

\subsection*{Das Beste kommt zum Schluss:}

\subsubsection*{Jedes Quantensystem hat sein eigenes Q.}
Ja, auch \textbf{DU}. Dein Geist? Quantensystem. Deine Gedanken? Q-Wirkung. 
\textit{»Wünsche ans Universum«} war früher Esoterik – heute Physik.

\subsubsection*{Kollektives Bewusstsein steuert die Geschichte.}
Die Natur ist demokratisch. Die Summe aller Q-Quellen bestimmt den Lauf.

\subsubsection*{Aber: Die Mächtigen manipulieren genau das.}
Social Media = Aufmerksamkeitssteuerung = Q-Manipulation. 
\textit{»Du glaubst, du willst das – aber du bist nur Erfüllungsgehilfe.«}

\subsection*{Die gute Nachricht:}
Du kannst dein Q \textbf{bewusst ausrichten}. 
\begin{center}
    \large\textbf{Aufmerksamkeit zurückholen. Fokus. Klarheit.}
\end{center}
Denn manipulierbar ist nur das Q, das nicht weiß, dass es eine Quelle ist.

\vspace{0.3cm}
\hrule
\vspace{0.3cm}

\begin{center}
    {\Large\bfseries WDBT+ – nicht nur eine Theorie, ein Lebensgefühl.}
    
    \vspace{0.2cm}
    \small
    \textit{»Und das Quantenpotential ward Fleisch und wohnte unter uns – nicht-lokal, versteht sich.«}
    
    \vspace{0.3cm}
    \textbf{Komm in die Community!} \\
    (Treffen jeden 3. Vollmond im fraktalen Wald bei D=2,71)
\end{center}

\end{document}